\chapter{Conclusion 結論與未來工作}

\section{研究總結}

本研究針對搭乘公車用餐者「餐廳是否可用」與「公車是否搭得上」必須同時成立的需求，提出以「餐廳＋公車搭乘方案」為決策單位的混合推薦方法。方法的核心是三項分離：以向量檢索處理語意召回、以知識圖譜提供有來源的結構化證據、以程式化的三值查核判定營業時間與搭乘條件。語言模型僅位於入口（需求整理）與出口（依判定紀錄的說明），營業時間、預算、班次與轉乘銜接的判斷皆交由程式與證據完成；使用者出發後，系統持續監控所選班次，延誤時改選附近餐廳並重新規劃。


在評估方面，本研究規劃了關鍵字、向量與混合召回及語言模型直接排名等對照組，並以連續輪詢、紀錄重播與延誤注入驗證行程中監控，同時報告違反率與覆蓋率，以避免以拒答換取低錯誤。本研究不要求每個模組都呈現正向效果；若向量召回或圖譜在小資料集上沒有實質增益，該結果同樣是有效發現，並可據以縮減系統複雜度。


\section{目前進度}

目前完成研究問題與假設的界定、文獻核對、系統設計與評估規劃，尚無實驗進度與結果。


\section{研究階段與驗收}

各階段為計畫順序，尚未指定經團隊確認的完成日期。


\begin{enumerate}
\item M1 規格凍結：完成需求條件、三值判定標準、引用核實與評估協定。
\item M2 餐廳資料：蒐集影片逐字稿與留言，萃取屬性並以官方 Places API 取得 Place ID，核對營業時間與價格。
\item M3 公車資料介接：實測旅運規劃與預估到站（N1）端點，確認回傳內容、更新間隔與覆蓋範圍；若無法取得所需資料，則縮減承諾並報告阻礙。
\item M4 系統與先導：完成九個模組的整合，以 50 筆查詢進行先導，確定排序指標與樣本數。
\item M5 凍結評估：執行對照組、連續輪詢重播、延誤注入與現場驗證，填入實際資料規模、結果與信賴區間。
\end{enumerate}

\section{預期成果}

預期交付可解釋的「餐廳＋公車搭乘方案」推薦原型、可追溯的資料與評估協定，以及各模組有效、無效或成本過高的實證依據，包含條件分類與三值判定標準、版本化的餐廳與公車資料、推薦與監控演算法、對照組比較、失效案例與限制分析。目前尚無結果表；未來報告將依實際測量填寫，不會把「預期提升」改寫成「實驗顯示」。


\section{未來工作}

\begin{itemize}
\item 多元交通方式：延伸到「開車或公車」的配對，與其他大眾運輸或自行車。
\item 預估誤差校正：累積足夠的預估到站與實際到站資料後，對預估誤差進行校正。
\end{itemize}
